\section{Постановка задачи}
\label{sec:Chapter1} \index{Chapter1}

% -----------------------------------------------------------------------
% 1.1 РЕКОМЕНДАТЕЛЬНАЯ СИСТЕМА КАК REPEATED LEARNING PROCESS
% -----------------------------------------------------------------------

\subsection{Рекомендательная система как процесс с обратной связью}
\label{subsec:rl_process}

Зафиксируем конечное множество пользователей $\mathcal{U}$ и конечное
множество объектов (айтемов) $\mathcal{I}$, $|\mathcal{U}|=N$,
$|\mathcal{I}|=M$. На каждом дискретном шаге времени $t=0,1,2,\ldots$
рекомендательная система выполняет три операции:

\begin{enumerate}
    \item \textbf{Рекомендация.} Для каждого пользователя $u \in \mathcal{U}$
          система выбирает подмножество объектов $S^{(t)}_u \subseteq
          \mathcal{I}$ для показа в соответствии с текущей политикой
          $\pi^{(t)}$.

    \item \textbf{Пользовательский отклик.} Пользователь $u$ взаимодействует
          с показанными объектами; формируется множество наблюдений
          $\mathcal{D}^{(t)}_u = \{(u, i, r^{(t)}_{ui}) : i \in S^{(t)}_u\}$,
          где $r^{(t)}_{ui} \in \mathbb{R}$ фиксирует реакцию (оценку,
          клик, время просмотра и т.\,д.).

    \item \textbf{Переобучение.} Параметры модели обновляются по накопленным
          данным $\mathcal{D}^{(t)} = \bigcup_{u} \mathcal{D}^{(t)}_u$, порождая
          политику $\pi^{(t+1)}$.
\end{enumerate}

Принципиальное отличие этого цикла от стандартной схемы обучения с учителем
состоит в том, что $\mathcal{D}^{(t)}$ есть функционал от политики $\pi^{(t)}$:
объекты, не попавшие в выдачу $S^{(t)}_u$, не порождают обучающих примеров
для шага $t{+}1$. Следовательно, задача оценки модели по наблюдаемым данным
принципиально смещена~--- это явление называется \textit{selection bias},
индуцированным самой рекомендательной системой.

\subsection{Пространство состояний и оператор переобучения}
\label{subsec:state_space}

Для описания долгосрочной динамики введём пространство состояний.
Обозначим через $\mathcal{P}(\mathcal{U})$ множество вероятностных мер
на $\mathcal{U}$; аналогично $\mathcal{P}(\mathcal{I})$~--- вероятностные
меры на $\mathcal{I}$. Состояние системы в момент $t$ задаётся парой
$(\mu^{(t)}, \nu^{(t)}) \in \mathcal{P}(\mathcal{U}) \times
\mathcal{P}(\mathcal{I})$, где $\mu^{(t)}$ описывает активность
пользователей, а $\nu^{(t)}$~--- распределение популярности объектов
в обучающей выборке.

\begin{definition}[Оператор переобучения]
\label{def:retraining_op}
\textit{Оператором переобучения} называется отображение
\[
    \mathcal{F} \colon \mathcal{P}(\mathcal{U}) \times
    \mathcal{P}(\mathcal{I}) \;\longrightarrow\;
    \Pi,
\]
сопоставляющее распределению данных $(\mu, \nu)$ политику рекомендаций
$\pi = \mathcal{F}(\mu, \nu)$, получаемую минимизацией функционала потерь
по обучающей выборке с распределением $(\mu, \nu)$.
\end{definition}

Совместно с оператором пользовательского отклика
$\mathcal{G} \colon \Pi \to \mathcal{P}(\mathcal{U}) \times \mathcal{P}(\mathcal{I})$
(преобразующим политику в распределение наблюдаемых реакций) оператор
переобучения задаёт дискретную динамическую систему
\[
    (\mu^{(t+1)},\, \nu^{(t+1)}) \;=\; \mathcal{G}\!\bigl(\mathcal{F}(\mu^{(t)},\, \nu^{(t)})\bigr).
\]

% -----------------------------------------------------------------------
% 1.2 КЛАСТЕРНАЯ СТРУКТУРА И ЛОКАЛЬНОЕ РАЗНООБРАЗИЕ
% -----------------------------------------------------------------------

\subsection{Кластерная структура аудитории}
\label{subsec:cluster_structure}

Предполагается, что пользователи образуют $K$ кластеров
$\mathcal{C}_1, \ldots, \mathcal{C}_K$ (например, по сегментам вкусовых
предпочтений). Пусть $\mu_k^{(t)} = \mu^{(t)}(\mathcal{C}_k)$~--- доля
кластера $k$ в момент $t$, и пусть $\overline{x}_k^{(t)}$~---
центроид (вектор средних предпочтений) кластера $k$, а
$\Sigma_k^{(t)}$~--- ковариационная матрица пользовательских
предпочтений внутри кластера $k$.

\begin{definition}[Внутрикластерное разнообразие]
\label{def:diversity}
\textit{Внутрикластерным разнообразием} в момент $t$ называется
суммарный след ковариационных матриц:
\[
    D^{(t)} \;=\; \sum_{k=1}^{K} \mu_k^{(t)} \,\operatorname{tr}\!\bigl(\Sigma_k^{(t)}\bigr).
\]
\end{definition}

\begin{definition}[Эхо-камера]
\label{def:echo_chamber}
Говорим, что система \textit{достигла режима эхо-камеры} к моменту $T$,
если для каждого кластера $k$ распределение рекомендаций в $\mathcal{C}_k$
сосредоточено в $\varepsilon$-окрестности единственного объекта
$i^*_k \in \mathcal{I}$:
\[
    \pi^{(T)}\bigl(S^{(T)}_u \ni i \bigr) < \varepsilon
    \quad \forall\, u \in \mathcal{C}_k,\; i \neq i^*_k.
\]
В пределе $\varepsilon \to 0$ это означает сходимость распределения
рекомендаций к $K$-модальной сумме дельта-мер:
\[
    \pi^{(\infty)} \;\xrightarrow{w}\; \sum_{k=1}^{K} \mu_k \,\delta_{i^*_k}.
\]
\end{definition}

% -----------------------------------------------------------------------
% 1.3 ФОРМУЛИРОВКА ИССЛЕДОВАТЕЛЬСКОГО ВОПРОСА
% -----------------------------------------------------------------------

\subsection{Исследовательский вопрос}
\label{subsec:research_question}

Таким образом, основной исследовательский вопрос формулируется следующим
образом.

\begin{quote}
\textit{При каких условиях на операторы $\mathcal{F}$ и $\mathcal{G}$,
кластерную структуру $\{\mathcal{C}_k\}$ и начальное состояние
$(\mu^{(0)}, \nu^{(0)})$ динамика системы приводит к обращению
внутрикластерного разнообразия $D^{(t)} \to 0$ при $t \to \infty$,
то есть к режиму $K$-модальной эхо-камеры?}
\end{quote}

Дополнительный вопрос касается управляемости: при каких значениях
параметра exploration $\alpha \in [0,1]$ и регуляризационного
коэффициента $\lambda > 0$ разнообразие $D^{(t)}$ сохраняется
ненулевым при $t \to \infty$?

Ответы на эти вопросы получены в главах~\ref{sec:Chapter2}
и~\ref{sec:Chapter3}.
